% ============================================================================
% ESQUEMA PROPOSTO (banca, sugestão — não mergeado): o programa experimental
% como MAPA ENCADEADO — o que a Tabela tab:metodo-experimentos lista, esta
% figura liga: qual experimento alimenta qual, o que é condicional, e onde
% cada resposta é lida. Complementa (não substitui) a tabela.
%
% Uso no Cap. 3 (dentro de um ambiente figure):
%   % ============================================================================
% ESQUEMA PROPOSTO (banca, sugestão — não mergeado): o programa experimental
% como MAPA ENCADEADO — o que a Tabela tab:metodo-experimentos lista, esta
% figura liga: qual experimento alimenta qual, o que é condicional, e onde
% cada resposta é lida. Complementa (não substitui) a tabela.
%
% Uso no Cap. 3 (dentro de um ambiente figure):
%   % ============================================================================
% ESQUEMA PROPOSTO (banca, sugestão — não mergeado): o programa experimental
% como MAPA ENCADEADO — o que a Tabela tab:metodo-experimentos lista, esta
% figura liga: qual experimento alimenta qual, o que é condicional, e onde
% cada resposta é lida. Complementa (não substitui) a tabela.
%
% Uso no Cap. 3 (dentro de um ambiente figure):
%   % ============================================================================
% ESQUEMA PROPOSTO (banca, sugestão — não mergeado): o programa experimental
% como MAPA ENCADEADO — o que a Tabela tab:metodo-experimentos lista, esta
% figura liga: qual experimento alimenta qual, o que é condicional, e onde
% cada resposta é lida. Complementa (não substitui) a tabela.
%
% Uso no Cap. 3 (dentro de um ambiente figure):
%   \input{3-metodo/esquemas-propostos/esq-mapa-experimental}
% Uso standalone (pré-visualização):
%   \documentclass[tikz,border=6pt]{standalone}
%   \usetikzlibrary{arrows.meta, positioning}
%   \begin{document}\input{esq-mapa-experimental.tex}\end{document}
%   (no modo standalone, troque as \ref{...} por texto fixo ou defina-as)
% Requer apenas arrows.meta e positioning (já em packages.tex). Legível em P&B:
% pilar = cinza escuro; experimento = branco; condicional = tracejado;
% destino de leitura = cinza claro.
% FREEZE respeitado: nenhum número; os identificadores E0..E6/E3 são os da
% tabela (controle interno, permitidos no Cap. 3).
% Compilado e inspecionado visualmente (pdflatex + render PNG) — ver
% NOTA-esquemas.md, iterações 4-5 do loop.
% ============================================================================
\begin{tikzpicture}[
    x=1cm, y=1cm,
    exp/.style={draw, rounded corners=2pt, align=center, text width=36mm,
                minimum height=11mm, inner sep=3pt, font=\footnotesize},
    pilar/.style={draw, rounded corners=2pt, align=center, text width=28mm,
                  minimum height=11mm, inner sep=3pt, font=\footnotesize\bfseries,
                  fill=gray!25},
    saida/.style={exp, fill=gray!8},
    cond/.style={exp, dashed},
    seta/.style={-{Stealth[length=2.4mm]}, thick},
    dep/.style={-{Stealth[length=2.4mm]}, thick, dashed},
    rot/.style={font=\scriptsize, align=center, inner sep=2pt}
]
    % ---------- linha 1: pilares 1-2 (Cap. 4) -----------------------------
    \node[pilar] (p1) at (0,0)    {conjunto inicial\\(partida a frio)};
    \node[exp]   (P1) at (4.8,0)
        {sensibilidade do $L_0$, envelope do AG, DRI-SL e reexecução independente};
    \node[saida] (cap4) at (9.8,0)
        {Capítulo~\ref{ch:resultados-l0}: composição e construção do conjunto inicial};
    \draw[seta] (p1) -- (P1);
    \draw[seta] (P1) -- (cap4);

    % ---------- linha 2: pilar 3, o oráculo (Cap. 5, primeira metade) ------
    \node[pilar] (p3) at (0,-2.6) {oráculo LLM};
    \node[exp]   (E0) at (4.8,-2.6)
        {\textbf{E0/E0-P}: avaliação fatorial dos oráculos e do \textit{prompt} (RQ1--RQ4)};
    \node[saida] (cap5o) at (9.8,-2.6)
        {Capítulo~\ref{ch:resultados-falco}: oráculos, custo e instrumento};
    \draw[seta] (p3) -- (E0);
    \draw[seta] (E0) -- (cap5o);

    % ---------- linha 3: apoios e o condicional ----------------------------
    % a condição do E4 vai no texto da própria caixa (rótulo sobre a seta
    % tracejada colidia com a linha — achado da inspeção visual)
    \node[pilar] (p4) at (0,-5.2) {framework\\integrado};
    \node[exp]   (E1) at (4.8,-5.2)
        {\textbf{E1/E2}: estratégias sob oráculo perfeito; épocas do BERTimbau (interno)};
    \node[cond]  (E4) at (9.8,-5.2)
        {\textbf{E4} (condicional --- o gate do E0 decide se ocorre): robustez do aprendizado ao ruído do oráculo};
    \draw[seta] (p4) -- (E1);
    \draw[dep]  (E0.south east) -- (E4.north west);

    % ---------- linha 4: execução e escala ---------------------------------
    \node[exp] (E5) at (4.8,-7.8)
        {\textbf{E5}: ciclo real de rotulagem com LLM e critério de parada};
    \node[exp] (E6) at (9.8,-7.8)
        {\textbf{E6}: seletores em escala populacional e viés de autoavaliação};

    % ---------- linha 5: a síntese (E3) ------------------------------------
    \node[exp] (E3) at (4.8,-10.6)
        {\textbf{E3}: avaliação da hipótese central com o classificador forte};
    \node[saida] (verd) at (9.8,-10.6)
        {veredito da hipótese (Seção~\ref{sec:res-e3p})};
    \draw[seta] (E5) -- node[rot, left=1pt, text width=26mm]
        {o conjunto coletado no ciclo real é um dos braços} (E3);
    % anchor=north west pendura o rótulo abaixo-à-direita do ponto,
    % inteiro do lado livre da diagonal (a linha cortava a 1ª linha)
    \draw[seta] (E6.south) -- node[rot, anchor=north west, pos=0.42, xshift=1mm,
        text width=34mm] {o padrão dos seletores é reavaliado no
        \textit{transformer}} (E3.north east);
    \draw[seta] (E3) -- (verd);

    \node[rot, text width=42mm, anchor=north east] (recebe) at (2.2,-10.0)
        {o E3 recebe ainda: a configuração de oráculos (do E0, via gate) e o
         tamanho de lote da Fase~2 (do E1)};
    \draw[dep] (recebe.east) -- (E3.west);
\end{tikzpicture}

% Uso standalone (pré-visualização):
%   \documentclass[tikz,border=6pt]{standalone}
%   \usetikzlibrary{arrows.meta, positioning}
%   \begin{document}% ============================================================================
% ESQUEMA PROPOSTO (banca, sugestão — não mergeado): o programa experimental
% como MAPA ENCADEADO — o que a Tabela tab:metodo-experimentos lista, esta
% figura liga: qual experimento alimenta qual, o que é condicional, e onde
% cada resposta é lida. Complementa (não substitui) a tabela.
%
% Uso no Cap. 3 (dentro de um ambiente figure):
%   \input{3-metodo/esquemas-propostos/esq-mapa-experimental}
% Uso standalone (pré-visualização):
%   \documentclass[tikz,border=6pt]{standalone}
%   \usetikzlibrary{arrows.meta, positioning}
%   \begin{document}\input{esq-mapa-experimental.tex}\end{document}
%   (no modo standalone, troque as \ref{...} por texto fixo ou defina-as)
% Requer apenas arrows.meta e positioning (já em packages.tex). Legível em P&B:
% pilar = cinza escuro; experimento = branco; condicional = tracejado;
% destino de leitura = cinza claro.
% FREEZE respeitado: nenhum número; os identificadores E0..E6/E3 são os da
% tabela (controle interno, permitidos no Cap. 3).
% Compilado e inspecionado visualmente (pdflatex + render PNG) — ver
% NOTA-esquemas.md, iterações 4-5 do loop.
% ============================================================================
\begin{tikzpicture}[
    x=1cm, y=1cm,
    exp/.style={draw, rounded corners=2pt, align=center, text width=36mm,
                minimum height=11mm, inner sep=3pt, font=\footnotesize},
    pilar/.style={draw, rounded corners=2pt, align=center, text width=28mm,
                  minimum height=11mm, inner sep=3pt, font=\footnotesize\bfseries,
                  fill=gray!25},
    saida/.style={exp, fill=gray!8},
    cond/.style={exp, dashed},
    seta/.style={-{Stealth[length=2.4mm]}, thick},
    dep/.style={-{Stealth[length=2.4mm]}, thick, dashed},
    rot/.style={font=\scriptsize, align=center, inner sep=2pt}
]
    % ---------- linha 1: pilares 1-2 (Cap. 4) -----------------------------
    \node[pilar] (p1) at (0,0)    {conjunto inicial\\(partida a frio)};
    \node[exp]   (P1) at (4.8,0)
        {sensibilidade do $L_0$, envelope do AG, DRI-SL e reexecução independente};
    \node[saida] (cap4) at (9.8,0)
        {Capítulo~\ref{ch:resultados-l0}: composição e construção do conjunto inicial};
    \draw[seta] (p1) -- (P1);
    \draw[seta] (P1) -- (cap4);

    % ---------- linha 2: pilar 3, o oráculo (Cap. 5, primeira metade) ------
    \node[pilar] (p3) at (0,-2.6) {oráculo LLM};
    \node[exp]   (E0) at (4.8,-2.6)
        {\textbf{E0/E0-P}: avaliação fatorial dos oráculos e do \textit{prompt} (RQ1--RQ4)};
    \node[saida] (cap5o) at (9.8,-2.6)
        {Capítulo~\ref{ch:resultados-falco}: oráculos, custo e instrumento};
    \draw[seta] (p3) -- (E0);
    \draw[seta] (E0) -- (cap5o);

    % ---------- linha 3: apoios e o condicional ----------------------------
    % a condição do E4 vai no texto da própria caixa (rótulo sobre a seta
    % tracejada colidia com a linha — achado da inspeção visual)
    \node[pilar] (p4) at (0,-5.2) {framework\\integrado};
    \node[exp]   (E1) at (4.8,-5.2)
        {\textbf{E1/E2}: estratégias sob oráculo perfeito; épocas do BERTimbau (interno)};
    \node[cond]  (E4) at (9.8,-5.2)
        {\textbf{E4} (condicional --- o gate do E0 decide se ocorre): robustez do aprendizado ao ruído do oráculo};
    \draw[seta] (p4) -- (E1);
    \draw[dep]  (E0.south east) -- (E4.north west);

    % ---------- linha 4: execução e escala ---------------------------------
    \node[exp] (E5) at (4.8,-7.8)
        {\textbf{E5}: ciclo real de rotulagem com LLM e critério de parada};
    \node[exp] (E6) at (9.8,-7.8)
        {\textbf{E6}: seletores em escala populacional e viés de autoavaliação};

    % ---------- linha 5: a síntese (E3) ------------------------------------
    \node[exp] (E3) at (4.8,-10.6)
        {\textbf{E3}: avaliação da hipótese central com o classificador forte};
    \node[saida] (verd) at (9.8,-10.6)
        {veredito da hipótese (Seção~\ref{sec:res-e3p})};
    \draw[seta] (E5) -- node[rot, left=1pt, text width=26mm]
        {o conjunto coletado no ciclo real é um dos braços} (E3);
    % anchor=north west pendura o rótulo abaixo-à-direita do ponto,
    % inteiro do lado livre da diagonal (a linha cortava a 1ª linha)
    \draw[seta] (E6.south) -- node[rot, anchor=north west, pos=0.42, xshift=1mm,
        text width=34mm] {o padrão dos seletores é reavaliado no
        \textit{transformer}} (E3.north east);
    \draw[seta] (E3) -- (verd);

    \node[rot, text width=42mm, anchor=north east] (recebe) at (2.2,-10.0)
        {o E3 recebe ainda: a configuração de oráculos (do E0, via gate) e o
         tamanho de lote da Fase~2 (do E1)};
    \draw[dep] (recebe.east) -- (E3.west);
\end{tikzpicture}
\end{document}
%   (no modo standalone, troque as \ref{...} por texto fixo ou defina-as)
% Requer apenas arrows.meta e positioning (já em packages.tex). Legível em P&B:
% pilar = cinza escuro; experimento = branco; condicional = tracejado;
% destino de leitura = cinza claro.
% FREEZE respeitado: nenhum número; os identificadores E0..E6/E3 são os da
% tabela (controle interno, permitidos no Cap. 3).
% Compilado e inspecionado visualmente (pdflatex + render PNG) — ver
% NOTA-esquemas.md, iterações 4-5 do loop.
% ============================================================================
\begin{tikzpicture}[
    x=1cm, y=1cm,
    exp/.style={draw, rounded corners=2pt, align=center, text width=36mm,
                minimum height=11mm, inner sep=3pt, font=\footnotesize},
    pilar/.style={draw, rounded corners=2pt, align=center, text width=28mm,
                  minimum height=11mm, inner sep=3pt, font=\footnotesize\bfseries,
                  fill=gray!25},
    saida/.style={exp, fill=gray!8},
    cond/.style={exp, dashed},
    seta/.style={-{Stealth[length=2.4mm]}, thick},
    dep/.style={-{Stealth[length=2.4mm]}, thick, dashed},
    rot/.style={font=\scriptsize, align=center, inner sep=2pt}
]
    % ---------- linha 1: pilares 1-2 (Cap. 4) -----------------------------
    \node[pilar] (p1) at (0,0)    {conjunto inicial\\(partida a frio)};
    \node[exp]   (P1) at (4.8,0)
        {sensibilidade do $L_0$, envelope do AG, DRI-SL e reexecução independente};
    \node[saida] (cap4) at (9.8,0)
        {Capítulo~\ref{ch:resultados-l0}: composição e construção do conjunto inicial};
    \draw[seta] (p1) -- (P1);
    \draw[seta] (P1) -- (cap4);

    % ---------- linha 2: pilar 3, o oráculo (Cap. 5, primeira metade) ------
    \node[pilar] (p3) at (0,-2.6) {oráculo LLM};
    \node[exp]   (E0) at (4.8,-2.6)
        {\textbf{E0/E0-P}: avaliação fatorial dos oráculos e do \textit{prompt} (RQ1--RQ4)};
    \node[saida] (cap5o) at (9.8,-2.6)
        {Capítulo~\ref{ch:resultados-falco}: oráculos, custo e instrumento};
    \draw[seta] (p3) -- (E0);
    \draw[seta] (E0) -- (cap5o);

    % ---------- linha 3: apoios e o condicional ----------------------------
    % a condição do E4 vai no texto da própria caixa (rótulo sobre a seta
    % tracejada colidia com a linha — achado da inspeção visual)
    \node[pilar] (p4) at (0,-5.2) {framework\\integrado};
    \node[exp]   (E1) at (4.8,-5.2)
        {\textbf{E1/E2}: estratégias sob oráculo perfeito; épocas do BERTimbau (interno)};
    \node[cond]  (E4) at (9.8,-5.2)
        {\textbf{E4} (condicional --- o gate do E0 decide se ocorre): robustez do aprendizado ao ruído do oráculo};
    \draw[seta] (p4) -- (E1);
    \draw[dep]  (E0.south east) -- (E4.north west);

    % ---------- linha 4: execução e escala ---------------------------------
    \node[exp] (E5) at (4.8,-7.8)
        {\textbf{E5}: ciclo real de rotulagem com LLM e critério de parada};
    \node[exp] (E6) at (9.8,-7.8)
        {\textbf{E6}: seletores em escala populacional e viés de autoavaliação};

    % ---------- linha 5: a síntese (E3) ------------------------------------
    \node[exp] (E3) at (4.8,-10.6)
        {\textbf{E3}: avaliação da hipótese central com o classificador forte};
    \node[saida] (verd) at (9.8,-10.6)
        {veredito da hipótese (Seção~\ref{sec:res-e3p})};
    \draw[seta] (E5) -- node[rot, left=1pt, text width=26mm]
        {o conjunto coletado no ciclo real é um dos braços} (E3);
    % anchor=north west pendura o rótulo abaixo-à-direita do ponto,
    % inteiro do lado livre da diagonal (a linha cortava a 1ª linha)
    \draw[seta] (E6.south) -- node[rot, anchor=north west, pos=0.42, xshift=1mm,
        text width=34mm] {o padrão dos seletores é reavaliado no
        \textit{transformer}} (E3.north east);
    \draw[seta] (E3) -- (verd);

    \node[rot, text width=42mm, anchor=north east] (recebe) at (2.2,-10.0)
        {o E3 recebe ainda: a configuração de oráculos (do E0, via gate) e o
         tamanho de lote da Fase~2 (do E1)};
    \draw[dep] (recebe.east) -- (E3.west);
\end{tikzpicture}

% Uso standalone (pré-visualização):
%   \documentclass[tikz,border=6pt]{standalone}
%   \usetikzlibrary{arrows.meta, positioning}
%   \begin{document}% ============================================================================
% ESQUEMA PROPOSTO (banca, sugestão — não mergeado): o programa experimental
% como MAPA ENCADEADO — o que a Tabela tab:metodo-experimentos lista, esta
% figura liga: qual experimento alimenta qual, o que é condicional, e onde
% cada resposta é lida. Complementa (não substitui) a tabela.
%
% Uso no Cap. 3 (dentro de um ambiente figure):
%   % ============================================================================
% ESQUEMA PROPOSTO (banca, sugestão — não mergeado): o programa experimental
% como MAPA ENCADEADO — o que a Tabela tab:metodo-experimentos lista, esta
% figura liga: qual experimento alimenta qual, o que é condicional, e onde
% cada resposta é lida. Complementa (não substitui) a tabela.
%
% Uso no Cap. 3 (dentro de um ambiente figure):
%   \input{3-metodo/esquemas-propostos/esq-mapa-experimental}
% Uso standalone (pré-visualização):
%   \documentclass[tikz,border=6pt]{standalone}
%   \usetikzlibrary{arrows.meta, positioning}
%   \begin{document}\input{esq-mapa-experimental.tex}\end{document}
%   (no modo standalone, troque as \ref{...} por texto fixo ou defina-as)
% Requer apenas arrows.meta e positioning (já em packages.tex). Legível em P&B:
% pilar = cinza escuro; experimento = branco; condicional = tracejado;
% destino de leitura = cinza claro.
% FREEZE respeitado: nenhum número; os identificadores E0..E6/E3 são os da
% tabela (controle interno, permitidos no Cap. 3).
% Compilado e inspecionado visualmente (pdflatex + render PNG) — ver
% NOTA-esquemas.md, iterações 4-5 do loop.
% ============================================================================
\begin{tikzpicture}[
    x=1cm, y=1cm,
    exp/.style={draw, rounded corners=2pt, align=center, text width=36mm,
                minimum height=11mm, inner sep=3pt, font=\footnotesize},
    pilar/.style={draw, rounded corners=2pt, align=center, text width=28mm,
                  minimum height=11mm, inner sep=3pt, font=\footnotesize\bfseries,
                  fill=gray!25},
    saida/.style={exp, fill=gray!8},
    cond/.style={exp, dashed},
    seta/.style={-{Stealth[length=2.4mm]}, thick},
    dep/.style={-{Stealth[length=2.4mm]}, thick, dashed},
    rot/.style={font=\scriptsize, align=center, inner sep=2pt}
]
    % ---------- linha 1: pilares 1-2 (Cap. 4) -----------------------------
    \node[pilar] (p1) at (0,0)    {conjunto inicial\\(partida a frio)};
    \node[exp]   (P1) at (4.8,0)
        {sensibilidade do $L_0$, envelope do AG, DRI-SL e reexecução independente};
    \node[saida] (cap4) at (9.8,0)
        {Capítulo~\ref{ch:resultados-l0}: composição e construção do conjunto inicial};
    \draw[seta] (p1) -- (P1);
    \draw[seta] (P1) -- (cap4);

    % ---------- linha 2: pilar 3, o oráculo (Cap. 5, primeira metade) ------
    \node[pilar] (p3) at (0,-2.6) {oráculo LLM};
    \node[exp]   (E0) at (4.8,-2.6)
        {\textbf{E0/E0-P}: avaliação fatorial dos oráculos e do \textit{prompt} (RQ1--RQ4)};
    \node[saida] (cap5o) at (9.8,-2.6)
        {Capítulo~\ref{ch:resultados-falco}: oráculos, custo e instrumento};
    \draw[seta] (p3) -- (E0);
    \draw[seta] (E0) -- (cap5o);

    % ---------- linha 3: apoios e o condicional ----------------------------
    % a condição do E4 vai no texto da própria caixa (rótulo sobre a seta
    % tracejada colidia com a linha — achado da inspeção visual)
    \node[pilar] (p4) at (0,-5.2) {framework\\integrado};
    \node[exp]   (E1) at (4.8,-5.2)
        {\textbf{E1/E2}: estratégias sob oráculo perfeito; épocas do BERTimbau (interno)};
    \node[cond]  (E4) at (9.8,-5.2)
        {\textbf{E4} (condicional --- o gate do E0 decide se ocorre): robustez do aprendizado ao ruído do oráculo};
    \draw[seta] (p4) -- (E1);
    \draw[dep]  (E0.south east) -- (E4.north west);

    % ---------- linha 4: execução e escala ---------------------------------
    \node[exp] (E5) at (4.8,-7.8)
        {\textbf{E5}: ciclo real de rotulagem com LLM e critério de parada};
    \node[exp] (E6) at (9.8,-7.8)
        {\textbf{E6}: seletores em escala populacional e viés de autoavaliação};

    % ---------- linha 5: a síntese (E3) ------------------------------------
    \node[exp] (E3) at (4.8,-10.6)
        {\textbf{E3}: avaliação da hipótese central com o classificador forte};
    \node[saida] (verd) at (9.8,-10.6)
        {veredito da hipótese (Seção~\ref{sec:res-e3p})};
    \draw[seta] (E5) -- node[rot, left=1pt, text width=26mm]
        {o conjunto coletado no ciclo real é um dos braços} (E3);
    % anchor=north west pendura o rótulo abaixo-à-direita do ponto,
    % inteiro do lado livre da diagonal (a linha cortava a 1ª linha)
    \draw[seta] (E6.south) -- node[rot, anchor=north west, pos=0.42, xshift=1mm,
        text width=34mm] {o padrão dos seletores é reavaliado no
        \textit{transformer}} (E3.north east);
    \draw[seta] (E3) -- (verd);

    \node[rot, text width=42mm, anchor=north east] (recebe) at (2.2,-10.0)
        {o E3 recebe ainda: a configuração de oráculos (do E0, via gate) e o
         tamanho de lote da Fase~2 (do E1)};
    \draw[dep] (recebe.east) -- (E3.west);
\end{tikzpicture}

% Uso standalone (pré-visualização):
%   \documentclass[tikz,border=6pt]{standalone}
%   \usetikzlibrary{arrows.meta, positioning}
%   \begin{document}% ============================================================================
% ESQUEMA PROPOSTO (banca, sugestão — não mergeado): o programa experimental
% como MAPA ENCADEADO — o que a Tabela tab:metodo-experimentos lista, esta
% figura liga: qual experimento alimenta qual, o que é condicional, e onde
% cada resposta é lida. Complementa (não substitui) a tabela.
%
% Uso no Cap. 3 (dentro de um ambiente figure):
%   \input{3-metodo/esquemas-propostos/esq-mapa-experimental}
% Uso standalone (pré-visualização):
%   \documentclass[tikz,border=6pt]{standalone}
%   \usetikzlibrary{arrows.meta, positioning}
%   \begin{document}\input{esq-mapa-experimental.tex}\end{document}
%   (no modo standalone, troque as \ref{...} por texto fixo ou defina-as)
% Requer apenas arrows.meta e positioning (já em packages.tex). Legível em P&B:
% pilar = cinza escuro; experimento = branco; condicional = tracejado;
% destino de leitura = cinza claro.
% FREEZE respeitado: nenhum número; os identificadores E0..E6/E3 são os da
% tabela (controle interno, permitidos no Cap. 3).
% Compilado e inspecionado visualmente (pdflatex + render PNG) — ver
% NOTA-esquemas.md, iterações 4-5 do loop.
% ============================================================================
\begin{tikzpicture}[
    x=1cm, y=1cm,
    exp/.style={draw, rounded corners=2pt, align=center, text width=36mm,
                minimum height=11mm, inner sep=3pt, font=\footnotesize},
    pilar/.style={draw, rounded corners=2pt, align=center, text width=28mm,
                  minimum height=11mm, inner sep=3pt, font=\footnotesize\bfseries,
                  fill=gray!25},
    saida/.style={exp, fill=gray!8},
    cond/.style={exp, dashed},
    seta/.style={-{Stealth[length=2.4mm]}, thick},
    dep/.style={-{Stealth[length=2.4mm]}, thick, dashed},
    rot/.style={font=\scriptsize, align=center, inner sep=2pt}
]
    % ---------- linha 1: pilares 1-2 (Cap. 4) -----------------------------
    \node[pilar] (p1) at (0,0)    {conjunto inicial\\(partida a frio)};
    \node[exp]   (P1) at (4.8,0)
        {sensibilidade do $L_0$, envelope do AG, DRI-SL e reexecução independente};
    \node[saida] (cap4) at (9.8,0)
        {Capítulo~\ref{ch:resultados-l0}: composição e construção do conjunto inicial};
    \draw[seta] (p1) -- (P1);
    \draw[seta] (P1) -- (cap4);

    % ---------- linha 2: pilar 3, o oráculo (Cap. 5, primeira metade) ------
    \node[pilar] (p3) at (0,-2.6) {oráculo LLM};
    \node[exp]   (E0) at (4.8,-2.6)
        {\textbf{E0/E0-P}: avaliação fatorial dos oráculos e do \textit{prompt} (RQ1--RQ4)};
    \node[saida] (cap5o) at (9.8,-2.6)
        {Capítulo~\ref{ch:resultados-falco}: oráculos, custo e instrumento};
    \draw[seta] (p3) -- (E0);
    \draw[seta] (E0) -- (cap5o);

    % ---------- linha 3: apoios e o condicional ----------------------------
    % a condição do E4 vai no texto da própria caixa (rótulo sobre a seta
    % tracejada colidia com a linha — achado da inspeção visual)
    \node[pilar] (p4) at (0,-5.2) {framework\\integrado};
    \node[exp]   (E1) at (4.8,-5.2)
        {\textbf{E1/E2}: estratégias sob oráculo perfeito; épocas do BERTimbau (interno)};
    \node[cond]  (E4) at (9.8,-5.2)
        {\textbf{E4} (condicional --- o gate do E0 decide se ocorre): robustez do aprendizado ao ruído do oráculo};
    \draw[seta] (p4) -- (E1);
    \draw[dep]  (E0.south east) -- (E4.north west);

    % ---------- linha 4: execução e escala ---------------------------------
    \node[exp] (E5) at (4.8,-7.8)
        {\textbf{E5}: ciclo real de rotulagem com LLM e critério de parada};
    \node[exp] (E6) at (9.8,-7.8)
        {\textbf{E6}: seletores em escala populacional e viés de autoavaliação};

    % ---------- linha 5: a síntese (E3) ------------------------------------
    \node[exp] (E3) at (4.8,-10.6)
        {\textbf{E3}: avaliação da hipótese central com o classificador forte};
    \node[saida] (verd) at (9.8,-10.6)
        {veredito da hipótese (Seção~\ref{sec:res-e3p})};
    \draw[seta] (E5) -- node[rot, left=1pt, text width=26mm]
        {o conjunto coletado no ciclo real é um dos braços} (E3);
    % anchor=north west pendura o rótulo abaixo-à-direita do ponto,
    % inteiro do lado livre da diagonal (a linha cortava a 1ª linha)
    \draw[seta] (E6.south) -- node[rot, anchor=north west, pos=0.42, xshift=1mm,
        text width=34mm] {o padrão dos seletores é reavaliado no
        \textit{transformer}} (E3.north east);
    \draw[seta] (E3) -- (verd);

    \node[rot, text width=42mm, anchor=north east] (recebe) at (2.2,-10.0)
        {o E3 recebe ainda: a configuração de oráculos (do E0, via gate) e o
         tamanho de lote da Fase~2 (do E1)};
    \draw[dep] (recebe.east) -- (E3.west);
\end{tikzpicture}
\end{document}
%   (no modo standalone, troque as \ref{...} por texto fixo ou defina-as)
% Requer apenas arrows.meta e positioning (já em packages.tex). Legível em P&B:
% pilar = cinza escuro; experimento = branco; condicional = tracejado;
% destino de leitura = cinza claro.
% FREEZE respeitado: nenhum número; os identificadores E0..E6/E3 são os da
% tabela (controle interno, permitidos no Cap. 3).
% Compilado e inspecionado visualmente (pdflatex + render PNG) — ver
% NOTA-esquemas.md, iterações 4-5 do loop.
% ============================================================================
\begin{tikzpicture}[
    x=1cm, y=1cm,
    exp/.style={draw, rounded corners=2pt, align=center, text width=36mm,
                minimum height=11mm, inner sep=3pt, font=\footnotesize},
    pilar/.style={draw, rounded corners=2pt, align=center, text width=28mm,
                  minimum height=11mm, inner sep=3pt, font=\footnotesize\bfseries,
                  fill=gray!25},
    saida/.style={exp, fill=gray!8},
    cond/.style={exp, dashed},
    seta/.style={-{Stealth[length=2.4mm]}, thick},
    dep/.style={-{Stealth[length=2.4mm]}, thick, dashed},
    rot/.style={font=\scriptsize, align=center, inner sep=2pt}
]
    % ---------- linha 1: pilares 1-2 (Cap. 4) -----------------------------
    \node[pilar] (p1) at (0,0)    {conjunto inicial\\(partida a frio)};
    \node[exp]   (P1) at (4.8,0)
        {sensibilidade do $L_0$, envelope do AG, DRI-SL e reexecução independente};
    \node[saida] (cap4) at (9.8,0)
        {Capítulo~\ref{ch:resultados-l0}: composição e construção do conjunto inicial};
    \draw[seta] (p1) -- (P1);
    \draw[seta] (P1) -- (cap4);

    % ---------- linha 2: pilar 3, o oráculo (Cap. 5, primeira metade) ------
    \node[pilar] (p3) at (0,-2.6) {oráculo LLM};
    \node[exp]   (E0) at (4.8,-2.6)
        {\textbf{E0/E0-P}: avaliação fatorial dos oráculos e do \textit{prompt} (RQ1--RQ4)};
    \node[saida] (cap5o) at (9.8,-2.6)
        {Capítulo~\ref{ch:resultados-falco}: oráculos, custo e instrumento};
    \draw[seta] (p3) -- (E0);
    \draw[seta] (E0) -- (cap5o);

    % ---------- linha 3: apoios e o condicional ----------------------------
    % a condição do E4 vai no texto da própria caixa (rótulo sobre a seta
    % tracejada colidia com a linha — achado da inspeção visual)
    \node[pilar] (p4) at (0,-5.2) {framework\\integrado};
    \node[exp]   (E1) at (4.8,-5.2)
        {\textbf{E1/E2}: estratégias sob oráculo perfeito; épocas do BERTimbau (interno)};
    \node[cond]  (E4) at (9.8,-5.2)
        {\textbf{E4} (condicional --- o gate do E0 decide se ocorre): robustez do aprendizado ao ruído do oráculo};
    \draw[seta] (p4) -- (E1);
    \draw[dep]  (E0.south east) -- (E4.north west);

    % ---------- linha 4: execução e escala ---------------------------------
    \node[exp] (E5) at (4.8,-7.8)
        {\textbf{E5}: ciclo real de rotulagem com LLM e critério de parada};
    \node[exp] (E6) at (9.8,-7.8)
        {\textbf{E6}: seletores em escala populacional e viés de autoavaliação};

    % ---------- linha 5: a síntese (E3) ------------------------------------
    \node[exp] (E3) at (4.8,-10.6)
        {\textbf{E3}: avaliação da hipótese central com o classificador forte};
    \node[saida] (verd) at (9.8,-10.6)
        {veredito da hipótese (Seção~\ref{sec:res-e3p})};
    \draw[seta] (E5) -- node[rot, left=1pt, text width=26mm]
        {o conjunto coletado no ciclo real é um dos braços} (E3);
    % anchor=north west pendura o rótulo abaixo-à-direita do ponto,
    % inteiro do lado livre da diagonal (a linha cortava a 1ª linha)
    \draw[seta] (E6.south) -- node[rot, anchor=north west, pos=0.42, xshift=1mm,
        text width=34mm] {o padrão dos seletores é reavaliado no
        \textit{transformer}} (E3.north east);
    \draw[seta] (E3) -- (verd);

    \node[rot, text width=42mm, anchor=north east] (recebe) at (2.2,-10.0)
        {o E3 recebe ainda: a configuração de oráculos (do E0, via gate) e o
         tamanho de lote da Fase~2 (do E1)};
    \draw[dep] (recebe.east) -- (E3.west);
\end{tikzpicture}
\end{document}
%   (no modo standalone, troque as \ref{...} por texto fixo ou defina-as)
% Requer apenas arrows.meta e positioning (já em packages.tex). Legível em P&B:
% pilar = cinza escuro; experimento = branco; condicional = tracejado;
% destino de leitura = cinza claro.
% FREEZE respeitado: nenhum número; os identificadores E0..E6/E3 são os da
% tabela (controle interno, permitidos no Cap. 3).
% Compilado e inspecionado visualmente (pdflatex + render PNG) — ver
% NOTA-esquemas.md, iterações 4-5 do loop.
% ============================================================================
\begin{tikzpicture}[
    x=1cm, y=1cm,
    exp/.style={draw, rounded corners=2pt, align=center, text width=36mm,
                minimum height=11mm, inner sep=3pt, font=\footnotesize},
    pilar/.style={draw, rounded corners=2pt, align=center, text width=28mm,
                  minimum height=11mm, inner sep=3pt, font=\footnotesize\bfseries,
                  fill=gray!25},
    saida/.style={exp, fill=gray!8},
    cond/.style={exp, dashed},
    seta/.style={-{Stealth[length=2.4mm]}, thick},
    dep/.style={-{Stealth[length=2.4mm]}, thick, dashed},
    rot/.style={font=\scriptsize, align=center, inner sep=2pt}
]
    % ---------- linha 1: pilares 1-2 (Cap. 4) -----------------------------
    \node[pilar] (p1) at (0,0)    {conjunto inicial\\(partida a frio)};
    \node[exp]   (P1) at (4.8,0)
        {sensibilidade do $L_0$, envelope do AG, DRI-SL e reexecução independente};
    \node[saida] (cap4) at (9.8,0)
        {Capítulo~\ref{ch:resultados-l0}: composição e construção do conjunto inicial};
    \draw[seta] (p1) -- (P1);
    \draw[seta] (P1) -- (cap4);

    % ---------- linha 2: pilar 3, o oráculo (Cap. 5, primeira metade) ------
    \node[pilar] (p3) at (0,-2.6) {oráculo LLM};
    \node[exp]   (E0) at (4.8,-2.6)
        {\textbf{E0/E0-P}: avaliação fatorial dos oráculos e do \textit{prompt} (RQ1--RQ4)};
    \node[saida] (cap5o) at (9.8,-2.6)
        {Capítulo~\ref{ch:resultados-falco}: oráculos, custo e instrumento};
    \draw[seta] (p3) -- (E0);
    \draw[seta] (E0) -- (cap5o);

    % ---------- linha 3: apoios e o condicional ----------------------------
    % a condição do E4 vai no texto da própria caixa (rótulo sobre a seta
    % tracejada colidia com a linha — achado da inspeção visual)
    \node[pilar] (p4) at (0,-5.2) {framework\\integrado};
    \node[exp]   (E1) at (4.8,-5.2)
        {\textbf{E1/E2}: estratégias sob oráculo perfeito; épocas do BERTimbau (interno)};
    \node[cond]  (E4) at (9.8,-5.2)
        {\textbf{E4} (condicional --- o gate do E0 decide se ocorre): robustez do aprendizado ao ruído do oráculo};
    \draw[seta] (p4) -- (E1);
    \draw[dep]  (E0.south east) -- (E4.north west);

    % ---------- linha 4: execução e escala ---------------------------------
    \node[exp] (E5) at (4.8,-7.8)
        {\textbf{E5}: ciclo real de rotulagem com LLM e critério de parada};
    \node[exp] (E6) at (9.8,-7.8)
        {\textbf{E6}: seletores em escala populacional e viés de autoavaliação};

    % ---------- linha 5: a síntese (E3) ------------------------------------
    \node[exp] (E3) at (4.8,-10.6)
        {\textbf{E3}: avaliação da hipótese central com o classificador forte};
    \node[saida] (verd) at (9.8,-10.6)
        {veredito da hipótese (Seção~\ref{sec:res-e3p})};
    \draw[seta] (E5) -- node[rot, left=1pt, text width=26mm]
        {o conjunto coletado no ciclo real é um dos braços} (E3);
    % anchor=north west pendura o rótulo abaixo-à-direita do ponto,
    % inteiro do lado livre da diagonal (a linha cortava a 1ª linha)
    \draw[seta] (E6.south) -- node[rot, anchor=north west, pos=0.42, xshift=1mm,
        text width=34mm] {o padrão dos seletores é reavaliado no
        \textit{transformer}} (E3.north east);
    \draw[seta] (E3) -- (verd);

    \node[rot, text width=42mm, anchor=north east] (recebe) at (2.2,-10.0)
        {o E3 recebe ainda: a configuração de oráculos (do E0, via gate) e o
         tamanho de lote da Fase~2 (do E1)};
    \draw[dep] (recebe.east) -- (E3.west);
\end{tikzpicture}

% Uso standalone (pré-visualização):
%   \documentclass[tikz,border=6pt]{standalone}
%   \usetikzlibrary{arrows.meta, positioning}
%   \begin{document}% ============================================================================
% ESQUEMA PROPOSTO (banca, sugestão — não mergeado): o programa experimental
% como MAPA ENCADEADO — o que a Tabela tab:metodo-experimentos lista, esta
% figura liga: qual experimento alimenta qual, o que é condicional, e onde
% cada resposta é lida. Complementa (não substitui) a tabela.
%
% Uso no Cap. 3 (dentro de um ambiente figure):
%   % ============================================================================
% ESQUEMA PROPOSTO (banca, sugestão — não mergeado): o programa experimental
% como MAPA ENCADEADO — o que a Tabela tab:metodo-experimentos lista, esta
% figura liga: qual experimento alimenta qual, o que é condicional, e onde
% cada resposta é lida. Complementa (não substitui) a tabela.
%
% Uso no Cap. 3 (dentro de um ambiente figure):
%   % ============================================================================
% ESQUEMA PROPOSTO (banca, sugestão — não mergeado): o programa experimental
% como MAPA ENCADEADO — o que a Tabela tab:metodo-experimentos lista, esta
% figura liga: qual experimento alimenta qual, o que é condicional, e onde
% cada resposta é lida. Complementa (não substitui) a tabela.
%
% Uso no Cap. 3 (dentro de um ambiente figure):
%   \input{3-metodo/esquemas-propostos/esq-mapa-experimental}
% Uso standalone (pré-visualização):
%   \documentclass[tikz,border=6pt]{standalone}
%   \usetikzlibrary{arrows.meta, positioning}
%   \begin{document}\input{esq-mapa-experimental.tex}\end{document}
%   (no modo standalone, troque as \ref{...} por texto fixo ou defina-as)
% Requer apenas arrows.meta e positioning (já em packages.tex). Legível em P&B:
% pilar = cinza escuro; experimento = branco; condicional = tracejado;
% destino de leitura = cinza claro.
% FREEZE respeitado: nenhum número; os identificadores E0..E6/E3 são os da
% tabela (controle interno, permitidos no Cap. 3).
% Compilado e inspecionado visualmente (pdflatex + render PNG) — ver
% NOTA-esquemas.md, iterações 4-5 do loop.
% ============================================================================
\begin{tikzpicture}[
    x=1cm, y=1cm,
    exp/.style={draw, rounded corners=2pt, align=center, text width=36mm,
                minimum height=11mm, inner sep=3pt, font=\footnotesize},
    pilar/.style={draw, rounded corners=2pt, align=center, text width=28mm,
                  minimum height=11mm, inner sep=3pt, font=\footnotesize\bfseries,
                  fill=gray!25},
    saida/.style={exp, fill=gray!8},
    cond/.style={exp, dashed},
    seta/.style={-{Stealth[length=2.4mm]}, thick},
    dep/.style={-{Stealth[length=2.4mm]}, thick, dashed},
    rot/.style={font=\scriptsize, align=center, inner sep=2pt}
]
    % ---------- linha 1: pilares 1-2 (Cap. 4) -----------------------------
    \node[pilar] (p1) at (0,0)    {conjunto inicial\\(partida a frio)};
    \node[exp]   (P1) at (4.8,0)
        {sensibilidade do $L_0$, envelope do AG, DRI-SL e reexecução independente};
    \node[saida] (cap4) at (9.8,0)
        {Capítulo~\ref{ch:resultados-l0}: composição e construção do conjunto inicial};
    \draw[seta] (p1) -- (P1);
    \draw[seta] (P1) -- (cap4);

    % ---------- linha 2: pilar 3, o oráculo (Cap. 5, primeira metade) ------
    \node[pilar] (p3) at (0,-2.6) {oráculo LLM};
    \node[exp]   (E0) at (4.8,-2.6)
        {\textbf{E0/E0-P}: avaliação fatorial dos oráculos e do \textit{prompt} (RQ1--RQ4)};
    \node[saida] (cap5o) at (9.8,-2.6)
        {Capítulo~\ref{ch:resultados-falco}: oráculos, custo e instrumento};
    \draw[seta] (p3) -- (E0);
    \draw[seta] (E0) -- (cap5o);

    % ---------- linha 3: apoios e o condicional ----------------------------
    % a condição do E4 vai no texto da própria caixa (rótulo sobre a seta
    % tracejada colidia com a linha — achado da inspeção visual)
    \node[pilar] (p4) at (0,-5.2) {framework\\integrado};
    \node[exp]   (E1) at (4.8,-5.2)
        {\textbf{E1/E2}: estratégias sob oráculo perfeito; épocas do BERTimbau (interno)};
    \node[cond]  (E4) at (9.8,-5.2)
        {\textbf{E4} (condicional --- o gate do E0 decide se ocorre): robustez do aprendizado ao ruído do oráculo};
    \draw[seta] (p4) -- (E1);
    \draw[dep]  (E0.south east) -- (E4.north west);

    % ---------- linha 4: execução e escala ---------------------------------
    \node[exp] (E5) at (4.8,-7.8)
        {\textbf{E5}: ciclo real de rotulagem com LLM e critério de parada};
    \node[exp] (E6) at (9.8,-7.8)
        {\textbf{E6}: seletores em escala populacional e viés de autoavaliação};

    % ---------- linha 5: a síntese (E3) ------------------------------------
    \node[exp] (E3) at (4.8,-10.6)
        {\textbf{E3}: avaliação da hipótese central com o classificador forte};
    \node[saida] (verd) at (9.8,-10.6)
        {veredito da hipótese (Seção~\ref{sec:res-e3p})};
    \draw[seta] (E5) -- node[rot, left=1pt, text width=26mm]
        {o conjunto coletado no ciclo real é um dos braços} (E3);
    % anchor=north west pendura o rótulo abaixo-à-direita do ponto,
    % inteiro do lado livre da diagonal (a linha cortava a 1ª linha)
    \draw[seta] (E6.south) -- node[rot, anchor=north west, pos=0.42, xshift=1mm,
        text width=34mm] {o padrão dos seletores é reavaliado no
        \textit{transformer}} (E3.north east);
    \draw[seta] (E3) -- (verd);

    \node[rot, text width=42mm, anchor=north east] (recebe) at (2.2,-10.0)
        {o E3 recebe ainda: a configuração de oráculos (do E0, via gate) e o
         tamanho de lote da Fase~2 (do E1)};
    \draw[dep] (recebe.east) -- (E3.west);
\end{tikzpicture}

% Uso standalone (pré-visualização):
%   \documentclass[tikz,border=6pt]{standalone}
%   \usetikzlibrary{arrows.meta, positioning}
%   \begin{document}% ============================================================================
% ESQUEMA PROPOSTO (banca, sugestão — não mergeado): o programa experimental
% como MAPA ENCADEADO — o que a Tabela tab:metodo-experimentos lista, esta
% figura liga: qual experimento alimenta qual, o que é condicional, e onde
% cada resposta é lida. Complementa (não substitui) a tabela.
%
% Uso no Cap. 3 (dentro de um ambiente figure):
%   \input{3-metodo/esquemas-propostos/esq-mapa-experimental}
% Uso standalone (pré-visualização):
%   \documentclass[tikz,border=6pt]{standalone}
%   \usetikzlibrary{arrows.meta, positioning}
%   \begin{document}\input{esq-mapa-experimental.tex}\end{document}
%   (no modo standalone, troque as \ref{...} por texto fixo ou defina-as)
% Requer apenas arrows.meta e positioning (já em packages.tex). Legível em P&B:
% pilar = cinza escuro; experimento = branco; condicional = tracejado;
% destino de leitura = cinza claro.
% FREEZE respeitado: nenhum número; os identificadores E0..E6/E3 são os da
% tabela (controle interno, permitidos no Cap. 3).
% Compilado e inspecionado visualmente (pdflatex + render PNG) — ver
% NOTA-esquemas.md, iterações 4-5 do loop.
% ============================================================================
\begin{tikzpicture}[
    x=1cm, y=1cm,
    exp/.style={draw, rounded corners=2pt, align=center, text width=36mm,
                minimum height=11mm, inner sep=3pt, font=\footnotesize},
    pilar/.style={draw, rounded corners=2pt, align=center, text width=28mm,
                  minimum height=11mm, inner sep=3pt, font=\footnotesize\bfseries,
                  fill=gray!25},
    saida/.style={exp, fill=gray!8},
    cond/.style={exp, dashed},
    seta/.style={-{Stealth[length=2.4mm]}, thick},
    dep/.style={-{Stealth[length=2.4mm]}, thick, dashed},
    rot/.style={font=\scriptsize, align=center, inner sep=2pt}
]
    % ---------- linha 1: pilares 1-2 (Cap. 4) -----------------------------
    \node[pilar] (p1) at (0,0)    {conjunto inicial\\(partida a frio)};
    \node[exp]   (P1) at (4.8,0)
        {sensibilidade do $L_0$, envelope do AG, DRI-SL e reexecução independente};
    \node[saida] (cap4) at (9.8,0)
        {Capítulo~\ref{ch:resultados-l0}: composição e construção do conjunto inicial};
    \draw[seta] (p1) -- (P1);
    \draw[seta] (P1) -- (cap4);

    % ---------- linha 2: pilar 3, o oráculo (Cap. 5, primeira metade) ------
    \node[pilar] (p3) at (0,-2.6) {oráculo LLM};
    \node[exp]   (E0) at (4.8,-2.6)
        {\textbf{E0/E0-P}: avaliação fatorial dos oráculos e do \textit{prompt} (RQ1--RQ4)};
    \node[saida] (cap5o) at (9.8,-2.6)
        {Capítulo~\ref{ch:resultados-falco}: oráculos, custo e instrumento};
    \draw[seta] (p3) -- (E0);
    \draw[seta] (E0) -- (cap5o);

    % ---------- linha 3: apoios e o condicional ----------------------------
    % a condição do E4 vai no texto da própria caixa (rótulo sobre a seta
    % tracejada colidia com a linha — achado da inspeção visual)
    \node[pilar] (p4) at (0,-5.2) {framework\\integrado};
    \node[exp]   (E1) at (4.8,-5.2)
        {\textbf{E1/E2}: estratégias sob oráculo perfeito; épocas do BERTimbau (interno)};
    \node[cond]  (E4) at (9.8,-5.2)
        {\textbf{E4} (condicional --- o gate do E0 decide se ocorre): robustez do aprendizado ao ruído do oráculo};
    \draw[seta] (p4) -- (E1);
    \draw[dep]  (E0.south east) -- (E4.north west);

    % ---------- linha 4: execução e escala ---------------------------------
    \node[exp] (E5) at (4.8,-7.8)
        {\textbf{E5}: ciclo real de rotulagem com LLM e critério de parada};
    \node[exp] (E6) at (9.8,-7.8)
        {\textbf{E6}: seletores em escala populacional e viés de autoavaliação};

    % ---------- linha 5: a síntese (E3) ------------------------------------
    \node[exp] (E3) at (4.8,-10.6)
        {\textbf{E3}: avaliação da hipótese central com o classificador forte};
    \node[saida] (verd) at (9.8,-10.6)
        {veredito da hipótese (Seção~\ref{sec:res-e3p})};
    \draw[seta] (E5) -- node[rot, left=1pt, text width=26mm]
        {o conjunto coletado no ciclo real é um dos braços} (E3);
    % anchor=north west pendura o rótulo abaixo-à-direita do ponto,
    % inteiro do lado livre da diagonal (a linha cortava a 1ª linha)
    \draw[seta] (E6.south) -- node[rot, anchor=north west, pos=0.42, xshift=1mm,
        text width=34mm] {o padrão dos seletores é reavaliado no
        \textit{transformer}} (E3.north east);
    \draw[seta] (E3) -- (verd);

    \node[rot, text width=42mm, anchor=north east] (recebe) at (2.2,-10.0)
        {o E3 recebe ainda: a configuração de oráculos (do E0, via gate) e o
         tamanho de lote da Fase~2 (do E1)};
    \draw[dep] (recebe.east) -- (E3.west);
\end{tikzpicture}
\end{document}
%   (no modo standalone, troque as \ref{...} por texto fixo ou defina-as)
% Requer apenas arrows.meta e positioning (já em packages.tex). Legível em P&B:
% pilar = cinza escuro; experimento = branco; condicional = tracejado;
% destino de leitura = cinza claro.
% FREEZE respeitado: nenhum número; os identificadores E0..E6/E3 são os da
% tabela (controle interno, permitidos no Cap. 3).
% Compilado e inspecionado visualmente (pdflatex + render PNG) — ver
% NOTA-esquemas.md, iterações 4-5 do loop.
% ============================================================================
\begin{tikzpicture}[
    x=1cm, y=1cm,
    exp/.style={draw, rounded corners=2pt, align=center, text width=36mm,
                minimum height=11mm, inner sep=3pt, font=\footnotesize},
    pilar/.style={draw, rounded corners=2pt, align=center, text width=28mm,
                  minimum height=11mm, inner sep=3pt, font=\footnotesize\bfseries,
                  fill=gray!25},
    saida/.style={exp, fill=gray!8},
    cond/.style={exp, dashed},
    seta/.style={-{Stealth[length=2.4mm]}, thick},
    dep/.style={-{Stealth[length=2.4mm]}, thick, dashed},
    rot/.style={font=\scriptsize, align=center, inner sep=2pt}
]
    % ---------- linha 1: pilares 1-2 (Cap. 4) -----------------------------
    \node[pilar] (p1) at (0,0)    {conjunto inicial\\(partida a frio)};
    \node[exp]   (P1) at (4.8,0)
        {sensibilidade do $L_0$, envelope do AG, DRI-SL e reexecução independente};
    \node[saida] (cap4) at (9.8,0)
        {Capítulo~\ref{ch:resultados-l0}: composição e construção do conjunto inicial};
    \draw[seta] (p1) -- (P1);
    \draw[seta] (P1) -- (cap4);

    % ---------- linha 2: pilar 3, o oráculo (Cap. 5, primeira metade) ------
    \node[pilar] (p3) at (0,-2.6) {oráculo LLM};
    \node[exp]   (E0) at (4.8,-2.6)
        {\textbf{E0/E0-P}: avaliação fatorial dos oráculos e do \textit{prompt} (RQ1--RQ4)};
    \node[saida] (cap5o) at (9.8,-2.6)
        {Capítulo~\ref{ch:resultados-falco}: oráculos, custo e instrumento};
    \draw[seta] (p3) -- (E0);
    \draw[seta] (E0) -- (cap5o);

    % ---------- linha 3: apoios e o condicional ----------------------------
    % a condição do E4 vai no texto da própria caixa (rótulo sobre a seta
    % tracejada colidia com a linha — achado da inspeção visual)
    \node[pilar] (p4) at (0,-5.2) {framework\\integrado};
    \node[exp]   (E1) at (4.8,-5.2)
        {\textbf{E1/E2}: estratégias sob oráculo perfeito; épocas do BERTimbau (interno)};
    \node[cond]  (E4) at (9.8,-5.2)
        {\textbf{E4} (condicional --- o gate do E0 decide se ocorre): robustez do aprendizado ao ruído do oráculo};
    \draw[seta] (p4) -- (E1);
    \draw[dep]  (E0.south east) -- (E4.north west);

    % ---------- linha 4: execução e escala ---------------------------------
    \node[exp] (E5) at (4.8,-7.8)
        {\textbf{E5}: ciclo real de rotulagem com LLM e critério de parada};
    \node[exp] (E6) at (9.8,-7.8)
        {\textbf{E6}: seletores em escala populacional e viés de autoavaliação};

    % ---------- linha 5: a síntese (E3) ------------------------------------
    \node[exp] (E3) at (4.8,-10.6)
        {\textbf{E3}: avaliação da hipótese central com o classificador forte};
    \node[saida] (verd) at (9.8,-10.6)
        {veredito da hipótese (Seção~\ref{sec:res-e3p})};
    \draw[seta] (E5) -- node[rot, left=1pt, text width=26mm]
        {o conjunto coletado no ciclo real é um dos braços} (E3);
    % anchor=north west pendura o rótulo abaixo-à-direita do ponto,
    % inteiro do lado livre da diagonal (a linha cortava a 1ª linha)
    \draw[seta] (E6.south) -- node[rot, anchor=north west, pos=0.42, xshift=1mm,
        text width=34mm] {o padrão dos seletores é reavaliado no
        \textit{transformer}} (E3.north east);
    \draw[seta] (E3) -- (verd);

    \node[rot, text width=42mm, anchor=north east] (recebe) at (2.2,-10.0)
        {o E3 recebe ainda: a configuração de oráculos (do E0, via gate) e o
         tamanho de lote da Fase~2 (do E1)};
    \draw[dep] (recebe.east) -- (E3.west);
\end{tikzpicture}

% Uso standalone (pré-visualização):
%   \documentclass[tikz,border=6pt]{standalone}
%   \usetikzlibrary{arrows.meta, positioning}
%   \begin{document}% ============================================================================
% ESQUEMA PROPOSTO (banca, sugestão — não mergeado): o programa experimental
% como MAPA ENCADEADO — o que a Tabela tab:metodo-experimentos lista, esta
% figura liga: qual experimento alimenta qual, o que é condicional, e onde
% cada resposta é lida. Complementa (não substitui) a tabela.
%
% Uso no Cap. 3 (dentro de um ambiente figure):
%   % ============================================================================
% ESQUEMA PROPOSTO (banca, sugestão — não mergeado): o programa experimental
% como MAPA ENCADEADO — o que a Tabela tab:metodo-experimentos lista, esta
% figura liga: qual experimento alimenta qual, o que é condicional, e onde
% cada resposta é lida. Complementa (não substitui) a tabela.
%
% Uso no Cap. 3 (dentro de um ambiente figure):
%   \input{3-metodo/esquemas-propostos/esq-mapa-experimental}
% Uso standalone (pré-visualização):
%   \documentclass[tikz,border=6pt]{standalone}
%   \usetikzlibrary{arrows.meta, positioning}
%   \begin{document}\input{esq-mapa-experimental.tex}\end{document}
%   (no modo standalone, troque as \ref{...} por texto fixo ou defina-as)
% Requer apenas arrows.meta e positioning (já em packages.tex). Legível em P&B:
% pilar = cinza escuro; experimento = branco; condicional = tracejado;
% destino de leitura = cinza claro.
% FREEZE respeitado: nenhum número; os identificadores E0..E6/E3 são os da
% tabela (controle interno, permitidos no Cap. 3).
% Compilado e inspecionado visualmente (pdflatex + render PNG) — ver
% NOTA-esquemas.md, iterações 4-5 do loop.
% ============================================================================
\begin{tikzpicture}[
    x=1cm, y=1cm,
    exp/.style={draw, rounded corners=2pt, align=center, text width=36mm,
                minimum height=11mm, inner sep=3pt, font=\footnotesize},
    pilar/.style={draw, rounded corners=2pt, align=center, text width=28mm,
                  minimum height=11mm, inner sep=3pt, font=\footnotesize\bfseries,
                  fill=gray!25},
    saida/.style={exp, fill=gray!8},
    cond/.style={exp, dashed},
    seta/.style={-{Stealth[length=2.4mm]}, thick},
    dep/.style={-{Stealth[length=2.4mm]}, thick, dashed},
    rot/.style={font=\scriptsize, align=center, inner sep=2pt}
]
    % ---------- linha 1: pilares 1-2 (Cap. 4) -----------------------------
    \node[pilar] (p1) at (0,0)    {conjunto inicial\\(partida a frio)};
    \node[exp]   (P1) at (4.8,0)
        {sensibilidade do $L_0$, envelope do AG, DRI-SL e reexecução independente};
    \node[saida] (cap4) at (9.8,0)
        {Capítulo~\ref{ch:resultados-l0}: composição e construção do conjunto inicial};
    \draw[seta] (p1) -- (P1);
    \draw[seta] (P1) -- (cap4);

    % ---------- linha 2: pilar 3, o oráculo (Cap. 5, primeira metade) ------
    \node[pilar] (p3) at (0,-2.6) {oráculo LLM};
    \node[exp]   (E0) at (4.8,-2.6)
        {\textbf{E0/E0-P}: avaliação fatorial dos oráculos e do \textit{prompt} (RQ1--RQ4)};
    \node[saida] (cap5o) at (9.8,-2.6)
        {Capítulo~\ref{ch:resultados-falco}: oráculos, custo e instrumento};
    \draw[seta] (p3) -- (E0);
    \draw[seta] (E0) -- (cap5o);

    % ---------- linha 3: apoios e o condicional ----------------------------
    % a condição do E4 vai no texto da própria caixa (rótulo sobre a seta
    % tracejada colidia com a linha — achado da inspeção visual)
    \node[pilar] (p4) at (0,-5.2) {framework\\integrado};
    \node[exp]   (E1) at (4.8,-5.2)
        {\textbf{E1/E2}: estratégias sob oráculo perfeito; épocas do BERTimbau (interno)};
    \node[cond]  (E4) at (9.8,-5.2)
        {\textbf{E4} (condicional --- o gate do E0 decide se ocorre): robustez do aprendizado ao ruído do oráculo};
    \draw[seta] (p4) -- (E1);
    \draw[dep]  (E0.south east) -- (E4.north west);

    % ---------- linha 4: execução e escala ---------------------------------
    \node[exp] (E5) at (4.8,-7.8)
        {\textbf{E5}: ciclo real de rotulagem com LLM e critério de parada};
    \node[exp] (E6) at (9.8,-7.8)
        {\textbf{E6}: seletores em escala populacional e viés de autoavaliação};

    % ---------- linha 5: a síntese (E3) ------------------------------------
    \node[exp] (E3) at (4.8,-10.6)
        {\textbf{E3}: avaliação da hipótese central com o classificador forte};
    \node[saida] (verd) at (9.8,-10.6)
        {veredito da hipótese (Seção~\ref{sec:res-e3p})};
    \draw[seta] (E5) -- node[rot, left=1pt, text width=26mm]
        {o conjunto coletado no ciclo real é um dos braços} (E3);
    % anchor=north west pendura o rótulo abaixo-à-direita do ponto,
    % inteiro do lado livre da diagonal (a linha cortava a 1ª linha)
    \draw[seta] (E6.south) -- node[rot, anchor=north west, pos=0.42, xshift=1mm,
        text width=34mm] {o padrão dos seletores é reavaliado no
        \textit{transformer}} (E3.north east);
    \draw[seta] (E3) -- (verd);

    \node[rot, text width=42mm, anchor=north east] (recebe) at (2.2,-10.0)
        {o E3 recebe ainda: a configuração de oráculos (do E0, via gate) e o
         tamanho de lote da Fase~2 (do E1)};
    \draw[dep] (recebe.east) -- (E3.west);
\end{tikzpicture}

% Uso standalone (pré-visualização):
%   \documentclass[tikz,border=6pt]{standalone}
%   \usetikzlibrary{arrows.meta, positioning}
%   \begin{document}% ============================================================================
% ESQUEMA PROPOSTO (banca, sugestão — não mergeado): o programa experimental
% como MAPA ENCADEADO — o que a Tabela tab:metodo-experimentos lista, esta
% figura liga: qual experimento alimenta qual, o que é condicional, e onde
% cada resposta é lida. Complementa (não substitui) a tabela.
%
% Uso no Cap. 3 (dentro de um ambiente figure):
%   \input{3-metodo/esquemas-propostos/esq-mapa-experimental}
% Uso standalone (pré-visualização):
%   \documentclass[tikz,border=6pt]{standalone}
%   \usetikzlibrary{arrows.meta, positioning}
%   \begin{document}\input{esq-mapa-experimental.tex}\end{document}
%   (no modo standalone, troque as \ref{...} por texto fixo ou defina-as)
% Requer apenas arrows.meta e positioning (já em packages.tex). Legível em P&B:
% pilar = cinza escuro; experimento = branco; condicional = tracejado;
% destino de leitura = cinza claro.
% FREEZE respeitado: nenhum número; os identificadores E0..E6/E3 são os da
% tabela (controle interno, permitidos no Cap. 3).
% Compilado e inspecionado visualmente (pdflatex + render PNG) — ver
% NOTA-esquemas.md, iterações 4-5 do loop.
% ============================================================================
\begin{tikzpicture}[
    x=1cm, y=1cm,
    exp/.style={draw, rounded corners=2pt, align=center, text width=36mm,
                minimum height=11mm, inner sep=3pt, font=\footnotesize},
    pilar/.style={draw, rounded corners=2pt, align=center, text width=28mm,
                  minimum height=11mm, inner sep=3pt, font=\footnotesize\bfseries,
                  fill=gray!25},
    saida/.style={exp, fill=gray!8},
    cond/.style={exp, dashed},
    seta/.style={-{Stealth[length=2.4mm]}, thick},
    dep/.style={-{Stealth[length=2.4mm]}, thick, dashed},
    rot/.style={font=\scriptsize, align=center, inner sep=2pt}
]
    % ---------- linha 1: pilares 1-2 (Cap. 4) -----------------------------
    \node[pilar] (p1) at (0,0)    {conjunto inicial\\(partida a frio)};
    \node[exp]   (P1) at (4.8,0)
        {sensibilidade do $L_0$, envelope do AG, DRI-SL e reexecução independente};
    \node[saida] (cap4) at (9.8,0)
        {Capítulo~\ref{ch:resultados-l0}: composição e construção do conjunto inicial};
    \draw[seta] (p1) -- (P1);
    \draw[seta] (P1) -- (cap4);

    % ---------- linha 2: pilar 3, o oráculo (Cap. 5, primeira metade) ------
    \node[pilar] (p3) at (0,-2.6) {oráculo LLM};
    \node[exp]   (E0) at (4.8,-2.6)
        {\textbf{E0/E0-P}: avaliação fatorial dos oráculos e do \textit{prompt} (RQ1--RQ4)};
    \node[saida] (cap5o) at (9.8,-2.6)
        {Capítulo~\ref{ch:resultados-falco}: oráculos, custo e instrumento};
    \draw[seta] (p3) -- (E0);
    \draw[seta] (E0) -- (cap5o);

    % ---------- linha 3: apoios e o condicional ----------------------------
    % a condição do E4 vai no texto da própria caixa (rótulo sobre a seta
    % tracejada colidia com a linha — achado da inspeção visual)
    \node[pilar] (p4) at (0,-5.2) {framework\\integrado};
    \node[exp]   (E1) at (4.8,-5.2)
        {\textbf{E1/E2}: estratégias sob oráculo perfeito; épocas do BERTimbau (interno)};
    \node[cond]  (E4) at (9.8,-5.2)
        {\textbf{E4} (condicional --- o gate do E0 decide se ocorre): robustez do aprendizado ao ruído do oráculo};
    \draw[seta] (p4) -- (E1);
    \draw[dep]  (E0.south east) -- (E4.north west);

    % ---------- linha 4: execução e escala ---------------------------------
    \node[exp] (E5) at (4.8,-7.8)
        {\textbf{E5}: ciclo real de rotulagem com LLM e critério de parada};
    \node[exp] (E6) at (9.8,-7.8)
        {\textbf{E6}: seletores em escala populacional e viés de autoavaliação};

    % ---------- linha 5: a síntese (E3) ------------------------------------
    \node[exp] (E3) at (4.8,-10.6)
        {\textbf{E3}: avaliação da hipótese central com o classificador forte};
    \node[saida] (verd) at (9.8,-10.6)
        {veredito da hipótese (Seção~\ref{sec:res-e3p})};
    \draw[seta] (E5) -- node[rot, left=1pt, text width=26mm]
        {o conjunto coletado no ciclo real é um dos braços} (E3);
    % anchor=north west pendura o rótulo abaixo-à-direita do ponto,
    % inteiro do lado livre da diagonal (a linha cortava a 1ª linha)
    \draw[seta] (E6.south) -- node[rot, anchor=north west, pos=0.42, xshift=1mm,
        text width=34mm] {o padrão dos seletores é reavaliado no
        \textit{transformer}} (E3.north east);
    \draw[seta] (E3) -- (verd);

    \node[rot, text width=42mm, anchor=north east] (recebe) at (2.2,-10.0)
        {o E3 recebe ainda: a configuração de oráculos (do E0, via gate) e o
         tamanho de lote da Fase~2 (do E1)};
    \draw[dep] (recebe.east) -- (E3.west);
\end{tikzpicture}
\end{document}
%   (no modo standalone, troque as \ref{...} por texto fixo ou defina-as)
% Requer apenas arrows.meta e positioning (já em packages.tex). Legível em P&B:
% pilar = cinza escuro; experimento = branco; condicional = tracejado;
% destino de leitura = cinza claro.
% FREEZE respeitado: nenhum número; os identificadores E0..E6/E3 são os da
% tabela (controle interno, permitidos no Cap. 3).
% Compilado e inspecionado visualmente (pdflatex + render PNG) — ver
% NOTA-esquemas.md, iterações 4-5 do loop.
% ============================================================================
\begin{tikzpicture}[
    x=1cm, y=1cm,
    exp/.style={draw, rounded corners=2pt, align=center, text width=36mm,
                minimum height=11mm, inner sep=3pt, font=\footnotesize},
    pilar/.style={draw, rounded corners=2pt, align=center, text width=28mm,
                  minimum height=11mm, inner sep=3pt, font=\footnotesize\bfseries,
                  fill=gray!25},
    saida/.style={exp, fill=gray!8},
    cond/.style={exp, dashed},
    seta/.style={-{Stealth[length=2.4mm]}, thick},
    dep/.style={-{Stealth[length=2.4mm]}, thick, dashed},
    rot/.style={font=\scriptsize, align=center, inner sep=2pt}
]
    % ---------- linha 1: pilares 1-2 (Cap. 4) -----------------------------
    \node[pilar] (p1) at (0,0)    {conjunto inicial\\(partida a frio)};
    \node[exp]   (P1) at (4.8,0)
        {sensibilidade do $L_0$, envelope do AG, DRI-SL e reexecução independente};
    \node[saida] (cap4) at (9.8,0)
        {Capítulo~\ref{ch:resultados-l0}: composição e construção do conjunto inicial};
    \draw[seta] (p1) -- (P1);
    \draw[seta] (P1) -- (cap4);

    % ---------- linha 2: pilar 3, o oráculo (Cap. 5, primeira metade) ------
    \node[pilar] (p3) at (0,-2.6) {oráculo LLM};
    \node[exp]   (E0) at (4.8,-2.6)
        {\textbf{E0/E0-P}: avaliação fatorial dos oráculos e do \textit{prompt} (RQ1--RQ4)};
    \node[saida] (cap5o) at (9.8,-2.6)
        {Capítulo~\ref{ch:resultados-falco}: oráculos, custo e instrumento};
    \draw[seta] (p3) -- (E0);
    \draw[seta] (E0) -- (cap5o);

    % ---------- linha 3: apoios e o condicional ----------------------------
    % a condição do E4 vai no texto da própria caixa (rótulo sobre a seta
    % tracejada colidia com a linha — achado da inspeção visual)
    \node[pilar] (p4) at (0,-5.2) {framework\\integrado};
    \node[exp]   (E1) at (4.8,-5.2)
        {\textbf{E1/E2}: estratégias sob oráculo perfeito; épocas do BERTimbau (interno)};
    \node[cond]  (E4) at (9.8,-5.2)
        {\textbf{E4} (condicional --- o gate do E0 decide se ocorre): robustez do aprendizado ao ruído do oráculo};
    \draw[seta] (p4) -- (E1);
    \draw[dep]  (E0.south east) -- (E4.north west);

    % ---------- linha 4: execução e escala ---------------------------------
    \node[exp] (E5) at (4.8,-7.8)
        {\textbf{E5}: ciclo real de rotulagem com LLM e critério de parada};
    \node[exp] (E6) at (9.8,-7.8)
        {\textbf{E6}: seletores em escala populacional e viés de autoavaliação};

    % ---------- linha 5: a síntese (E3) ------------------------------------
    \node[exp] (E3) at (4.8,-10.6)
        {\textbf{E3}: avaliação da hipótese central com o classificador forte};
    \node[saida] (verd) at (9.8,-10.6)
        {veredito da hipótese (Seção~\ref{sec:res-e3p})};
    \draw[seta] (E5) -- node[rot, left=1pt, text width=26mm]
        {o conjunto coletado no ciclo real é um dos braços} (E3);
    % anchor=north west pendura o rótulo abaixo-à-direita do ponto,
    % inteiro do lado livre da diagonal (a linha cortava a 1ª linha)
    \draw[seta] (E6.south) -- node[rot, anchor=north west, pos=0.42, xshift=1mm,
        text width=34mm] {o padrão dos seletores é reavaliado no
        \textit{transformer}} (E3.north east);
    \draw[seta] (E3) -- (verd);

    \node[rot, text width=42mm, anchor=north east] (recebe) at (2.2,-10.0)
        {o E3 recebe ainda: a configuração de oráculos (do E0, via gate) e o
         tamanho de lote da Fase~2 (do E1)};
    \draw[dep] (recebe.east) -- (E3.west);
\end{tikzpicture}
\end{document}
%   (no modo standalone, troque as \ref{...} por texto fixo ou defina-as)
% Requer apenas arrows.meta e positioning (já em packages.tex). Legível em P&B:
% pilar = cinza escuro; experimento = branco; condicional = tracejado;
% destino de leitura = cinza claro.
% FREEZE respeitado: nenhum número; os identificadores E0..E6/E3 são os da
% tabela (controle interno, permitidos no Cap. 3).
% Compilado e inspecionado visualmente (pdflatex + render PNG) — ver
% NOTA-esquemas.md, iterações 4-5 do loop.
% ============================================================================
\begin{tikzpicture}[
    x=1cm, y=1cm,
    exp/.style={draw, rounded corners=2pt, align=center, text width=36mm,
                minimum height=11mm, inner sep=3pt, font=\footnotesize},
    pilar/.style={draw, rounded corners=2pt, align=center, text width=28mm,
                  minimum height=11mm, inner sep=3pt, font=\footnotesize\bfseries,
                  fill=gray!25},
    saida/.style={exp, fill=gray!8},
    cond/.style={exp, dashed},
    seta/.style={-{Stealth[length=2.4mm]}, thick},
    dep/.style={-{Stealth[length=2.4mm]}, thick, dashed},
    rot/.style={font=\scriptsize, align=center, inner sep=2pt}
]
    % ---------- linha 1: pilares 1-2 (Cap. 4) -----------------------------
    \node[pilar] (p1) at (0,0)    {conjunto inicial\\(partida a frio)};
    \node[exp]   (P1) at (4.8,0)
        {sensibilidade do $L_0$, envelope do AG, DRI-SL e reexecução independente};
    \node[saida] (cap4) at (9.8,0)
        {Capítulo~\ref{ch:resultados-l0}: composição e construção do conjunto inicial};
    \draw[seta] (p1) -- (P1);
    \draw[seta] (P1) -- (cap4);

    % ---------- linha 2: pilar 3, o oráculo (Cap. 5, primeira metade) ------
    \node[pilar] (p3) at (0,-2.6) {oráculo LLM};
    \node[exp]   (E0) at (4.8,-2.6)
        {\textbf{E0/E0-P}: avaliação fatorial dos oráculos e do \textit{prompt} (RQ1--RQ4)};
    \node[saida] (cap5o) at (9.8,-2.6)
        {Capítulo~\ref{ch:resultados-falco}: oráculos, custo e instrumento};
    \draw[seta] (p3) -- (E0);
    \draw[seta] (E0) -- (cap5o);

    % ---------- linha 3: apoios e o condicional ----------------------------
    % a condição do E4 vai no texto da própria caixa (rótulo sobre a seta
    % tracejada colidia com a linha — achado da inspeção visual)
    \node[pilar] (p4) at (0,-5.2) {framework\\integrado};
    \node[exp]   (E1) at (4.8,-5.2)
        {\textbf{E1/E2}: estratégias sob oráculo perfeito; épocas do BERTimbau (interno)};
    \node[cond]  (E4) at (9.8,-5.2)
        {\textbf{E4} (condicional --- o gate do E0 decide se ocorre): robustez do aprendizado ao ruído do oráculo};
    \draw[seta] (p4) -- (E1);
    \draw[dep]  (E0.south east) -- (E4.north west);

    % ---------- linha 4: execução e escala ---------------------------------
    \node[exp] (E5) at (4.8,-7.8)
        {\textbf{E5}: ciclo real de rotulagem com LLM e critério de parada};
    \node[exp] (E6) at (9.8,-7.8)
        {\textbf{E6}: seletores em escala populacional e viés de autoavaliação};

    % ---------- linha 5: a síntese (E3) ------------------------------------
    \node[exp] (E3) at (4.8,-10.6)
        {\textbf{E3}: avaliação da hipótese central com o classificador forte};
    \node[saida] (verd) at (9.8,-10.6)
        {veredito da hipótese (Seção~\ref{sec:res-e3p})};
    \draw[seta] (E5) -- node[rot, left=1pt, text width=26mm]
        {o conjunto coletado no ciclo real é um dos braços} (E3);
    % anchor=north west pendura o rótulo abaixo-à-direita do ponto,
    % inteiro do lado livre da diagonal (a linha cortava a 1ª linha)
    \draw[seta] (E6.south) -- node[rot, anchor=north west, pos=0.42, xshift=1mm,
        text width=34mm] {o padrão dos seletores é reavaliado no
        \textit{transformer}} (E3.north east);
    \draw[seta] (E3) -- (verd);

    \node[rot, text width=42mm, anchor=north east] (recebe) at (2.2,-10.0)
        {o E3 recebe ainda: a configuração de oráculos (do E0, via gate) e o
         tamanho de lote da Fase~2 (do E1)};
    \draw[dep] (recebe.east) -- (E3.west);
\end{tikzpicture}
\end{document}
%   (no modo standalone, troque as \ref{...} por texto fixo ou defina-as)
% Requer apenas arrows.meta e positioning (já em packages.tex). Legível em P&B:
% pilar = cinza escuro; experimento = branco; condicional = tracejado;
% destino de leitura = cinza claro.
% FREEZE respeitado: nenhum número; os identificadores E0..E6/E3 são os da
% tabela (controle interno, permitidos no Cap. 3).
% Compilado e inspecionado visualmente (pdflatex + render PNG) — ver
% NOTA-esquemas.md, iterações 4-5 do loop.
% ============================================================================
\begin{tikzpicture}[
    x=1cm, y=1cm,
    exp/.style={draw, rounded corners=2pt, align=center, text width=36mm,
                minimum height=11mm, inner sep=3pt, font=\footnotesize},
    pilar/.style={draw, rounded corners=2pt, align=center, text width=28mm,
                  minimum height=11mm, inner sep=3pt, font=\footnotesize\bfseries,
                  fill=gray!25},
    saida/.style={exp, fill=gray!8},
    cond/.style={exp, dashed},
    seta/.style={-{Stealth[length=2.4mm]}, thick},
    dep/.style={-{Stealth[length=2.4mm]}, thick, dashed},
    rot/.style={font=\scriptsize, align=center, inner sep=2pt}
]
    % ---------- linha 1: pilares 1-2 (Cap. 4) -----------------------------
    \node[pilar] (p1) at (0,0)    {conjunto inicial\\(partida a frio)};
    \node[exp]   (P1) at (4.8,0)
        {sensibilidade do $L_0$, envelope do AG, DRI-SL e reexecução independente};
    \node[saida] (cap4) at (9.8,0)
        {Capítulo~\ref{ch:resultados-l0}: composição e construção do conjunto inicial};
    \draw[seta] (p1) -- (P1);
    \draw[seta] (P1) -- (cap4);

    % ---------- linha 2: pilar 3, o oráculo (Cap. 5, primeira metade) ------
    \node[pilar] (p3) at (0,-2.6) {oráculo LLM};
    \node[exp]   (E0) at (4.8,-2.6)
        {\textbf{E0/E0-P}: avaliação fatorial dos oráculos e do \textit{prompt} (RQ1--RQ4)};
    \node[saida] (cap5o) at (9.8,-2.6)
        {Capítulo~\ref{ch:resultados-falco}: oráculos, custo e instrumento};
    \draw[seta] (p3) -- (E0);
    \draw[seta] (E0) -- (cap5o);

    % ---------- linha 3: apoios e o condicional ----------------------------
    % a condição do E4 vai no texto da própria caixa (rótulo sobre a seta
    % tracejada colidia com a linha — achado da inspeção visual)
    \node[pilar] (p4) at (0,-5.2) {framework\\integrado};
    \node[exp]   (E1) at (4.8,-5.2)
        {\textbf{E1/E2}: estratégias sob oráculo perfeito; épocas do BERTimbau (interno)};
    \node[cond]  (E4) at (9.8,-5.2)
        {\textbf{E4} (condicional --- o gate do E0 decide se ocorre): robustez do aprendizado ao ruído do oráculo};
    \draw[seta] (p4) -- (E1);
    \draw[dep]  (E0.south east) -- (E4.north west);

    % ---------- linha 4: execução e escala ---------------------------------
    \node[exp] (E5) at (4.8,-7.8)
        {\textbf{E5}: ciclo real de rotulagem com LLM e critério de parada};
    \node[exp] (E6) at (9.8,-7.8)
        {\textbf{E6}: seletores em escala populacional e viés de autoavaliação};

    % ---------- linha 5: a síntese (E3) ------------------------------------
    \node[exp] (E3) at (4.8,-10.6)
        {\textbf{E3}: avaliação da hipótese central com o classificador forte};
    \node[saida] (verd) at (9.8,-10.6)
        {veredito da hipótese (Seção~\ref{sec:res-e3p})};
    \draw[seta] (E5) -- node[rot, left=1pt, text width=26mm]
        {o conjunto coletado no ciclo real é um dos braços} (E3);
    % anchor=north west pendura o rótulo abaixo-à-direita do ponto,
    % inteiro do lado livre da diagonal (a linha cortava a 1ª linha)
    \draw[seta] (E6.south) -- node[rot, anchor=north west, pos=0.42, xshift=1mm,
        text width=34mm] {o padrão dos seletores é reavaliado no
        \textit{transformer}} (E3.north east);
    \draw[seta] (E3) -- (verd);

    \node[rot, text width=42mm, anchor=north east] (recebe) at (2.2,-10.0)
        {o E3 recebe ainda: a configuração de oráculos (do E0, via gate) e o
         tamanho de lote da Fase~2 (do E1)};
    \draw[dep] (recebe.east) -- (E3.west);
\end{tikzpicture}
